% Part: incompleteness
% Chapter: arithmetization-syntax
% Section: coding-terms

\documentclass[../../../include/open-logic-section]{subfiles}

\begin{document}

\olfileid{inc}{art}{trm}
\olsection{પદોનું સંકેતબદ્ધીકરણ}

\begin{explain}
પદ એ ચિહ્નોના અનુક્રમનો એક ખાસ પ્રકાર
જ છે: પદોની રચનાના નિયમો મુજબ તે
સ્થિરો અને ચલોમાંથી આગમનાત્મક રીતે
બને છે. ચિહ્નોને સંકેતબદ્ધ કરવાની યોજના
અને સંખ્યાઓના અનુક્રમોને સંકેતબદ્ધ કરવાની
રીત વડે ચિહ્નોના અનુક્રમોને સંખ્યાઓ તરીકે
સંકેતબદ્ધ કરી શકાય છે. તેથી પદોને
ગડેલ-સંખ્યાઓ આપવી મુશ્કેલ નથી.
ખરો પડકાર તો એ બતાવવાનો છે કે
યોગ્ય રીતે રચાયેલા પદની ગડેલ-સંખ્યા
હોવાનો ગુણધર્મ સંગણનીય, ખરેખર આદિમ
પુનરાવર્તી છે.
\end{explain}

!!^{variable}s અને !!{constant}s સૌથી
સરળ પદો છે. $x$ આવા પદની ગડેલ-સંખ્યા
છે કે નહીં તે તપાસવું સહેલું છે:
કોઈ~$i$ માટે~$x$ એ~$\Gn{\Obj v_i}$
હોય ત્યારે $\fn{Var}(x)$ છે. બીજા
શબ્દોમાં, $x$ એ લંબાઈ~$1$નો અનુક્રમ છે
અને તેનો એકમાત્ર ઘટક~$(x)_0$ એ
કોઈ !!{variable}~$\Obj v_i$નો સંકેત
નંબર છે; એટલે કોઈ~$i$ માટે~$x$
એ~$\tuple{\tuple{1, i}}$ છે. એ જ રીતે,
કોઈ~$i$ માટે~$x$ એ~$\Gn{\Obj c_i}$
હોય ત્યારે~$\fn{Const}(x)$ છે. આ
બંને સંબંધો આદિમ પુનરાવર્તી છે, કારણ કે
આવો~$i$ હોય તો તે~$< x$ હોવો જ
જોઈએ:
\begin{align*}
  \fn{Var}(x) & \defiff \bexists{i<x}{x = \tuple{\tuple{1, i}}}\\
  \fn{Const}(x) & \defiff \bexists{i<x}{x = \tuple{\tuple{2, i}}}
\end{align*}

\begin{prop}
\ollabel{prop:term-primrec}
$x$ અનુક્રમે પદ અથવા બંધ પદની
ગડેલ-સંખ્યા હોય ત્યારે અને માત્ર ત્યારે
સત્ય થતા સંબંધો~$\fn{Term}(x)$ અને
$\fn{ClTerm}(x)$ આદિમ પુનરાવર્તી છે.
\end{prop}

\begin{proof}
ચિહ્નોનો અનુક્રમ~$s$ પદ છે ત્યારે અને
માત્ર ત્યારે, જ્યારે પદોનો એવો અનુક્રમ
$s_0$, \dots, $s_{k-1} = s$ હોય જે
પદોના રચના-નિયમો મુજબ~$s$ !!{constant}s અને
!!{variable}sમાંથી કેવી રીતે બન્યું તે નોંધે છે.
આવી સંભવિત રચના-શ્રેણી રચના-નિયમો
અનુસરે એ માટે દરેક~$i < k$ માટે
નીચેનામાંથી એક હોવું જોઈએ:
\begin{enumerate}
\item $s_i$ કોઈ !!a{variable}~$\Obj v_j$ છે, અથવા
\item $s_i$ કોઈ !!a{constant}~$\Obj c_j$ છે, અથવા
\item $s_i$ એ સ્થાન~$i$ પહેલાં આવતાં
  $n$ પદો~$t_1$, \dots, $t_n$માંથી
  $n$-સ્થાનિક !!{function}~$\Obj f^n_j$
  વાપરીને બનાવેલું છે.
\end{enumerate}
ગડેલ-સંખ્યાઓ પરનો અનુરૂપ સંબંધ આદિમ
પુનરાવર્તી છે તે બતાવવા આ શરતને
આદિમ પુનરાવર્તી રીતે—એટલે આદિમ
પુનરાવર્તી વિધેયો, સંબંધો અને સીમિત
પરિમાણન વડે—વ્યક્ત કરવી પડશે.

માનો કે~$y$ એ અનુક્રમ~$s_0$, \dots,
$s_{k-1}$નો સંકેત નંબર છે; એટલે
$y = \tuple{\Gn{s_0}, \dots, \Gn{s_{k-1}}}$.
દરેક~$i < k$ માટે નીચેની શરતો
સંતોષાય ત્યારે અને માત્ર ત્યારે તે
ગડેલ-સંખ્યા~$x$ ધરાવતા પદનો
રચના-શ્રેણી સંકેતબદ્ધ કરે છે:
\begin{enumerate}
\item $\fn{Var}((y)_i)$, અથવા
\item $\fn{Const}((y)_i)$, અથવા
\item કોઈ~$n$ અને
  નંબર~$z = \tuple{z_1, \dots, z_n}$
  હોય, જ્યાં દરેક~$z_l$ કોઈ~$i' < i$
  માટે~$(y)_{i'}$ જેટલો હોય, અને
\[
(y)_i = \Gn{\Obj f^n_j(} \concat \fn{flatten}(z) \concat \Gn{)},
\]
\end{enumerate}
અને વધુમાં~$(y)_{k-1} = x$ હોય.
(વિધેય~$\fn{flatten}(z)$ અનુક્રમ
$\tuple{\Gn{t_1}, \dots, \Gn{t_n}}$ને
$\Gn{t_1, \dots,
  t_n}$માં ફેરવે છે અને આદિમ પુનરાવર્તી છે.)

કિસ્સા~(3)માં સૂચકાંકો~$j$, $n$, પદો~$t_l$ની
ગડેલ-સંખ્યાઓ~$z_l$ અને
અનુક્રમ~$\tuple{z_1, \dots, z_n}$નો
સંકેત નંબર~$z$—આ બધાં~$y$ કરતાં
નાનાં છે. ઉપરનું~$k$ આપણે~$\len{y}$થી
બદલી શકીએ. તેથી ``$y$ એ ગડેલ-સંખ્યા~$x$
ધરાવતા પદની રચના-શ્રેણીનો સંકેત
નંબર છે'' એ શરતને અનુરૂપ સંબંધ
આદિમ પુનરાવર્તી છે તે સ્પષ્ટ થાય
એવી રીતે વ્યક્ત કરી શકીએ છીએ.

હવે~$y$ માટે આદિમ પુનરાવર્તી મર્યાદા
છે એટલું બતાવવાનું રહે છે. $x$ પદની
ગડેલ-સંખ્યા હોય, તો વધુમાં વધુ~$\len{x}$
પદોની રચના-શ્રેણી મળી શકે છે
(કારણ કે~$s$ની રચના-શ્રેણીનું દરેક પદ
$s$માં કોઈ સ્થાનથી શરૂ થવું જોઈએ,
અને બે ઉપપદો એક જ સ્થાનથી શરૂ
થઈ શકતાં નથી). $s$ના દરેક ઉપપદની
ગડેલ-સંખ્યા નિશ્ચિતપણે~$\le x$ છે.
તેથી હંમેશાં એવી રચના-શ્રેણી છે જેનો
સંકેત નંબર~$\le p_{k-1}^{k(x+1)}$ છે,
જ્યાં~$k=\len{x}$.

$\fn{ClTerm}$ માટે માત્ર
!!{variable}sવાળી શરત કાઢી નાખો.
\end{proof}

\begin{prob}
અનુક્રમ~$\tuple{\Gn{t_1}, \dots, \Gn{t_n}}$ને
$\Gn{t_1, \dots, t_n}$માં ફેરવતું
વિધેય~$\fn{flatten}(z)$ આદિમ
પુનરાવર્તી છે તે બતાવો.
\end{prob}

\begin{prop}\ollabel{prop:num-primrec}
  વિધેય~$\fn{num}(n) = \Gn{\num{n}}$ આદિમ પુનરાવર્તી છે.
\end{prop}

\begin{proof}
  $\fn{num}(n)$ને આદિમ પુનરાવર્તનથી વ્યાખ્યાયિત કરીએ:
  \begin{align*}
    \fn{num}(0) & = \Gn{\Obj 0}\\
    \fn{num}(n+1) & = \Gn{\prime(} \concat \fn{num}(n) \concat \Gn{)}.
  \end{align*}
\end{proof}

\end{document}
